\section*{Введение}
\addcontentsline{toc}{section}{Введение}

\begin{definition}
Алгебраическое (возможно, бесконечное) расширение $\Om/F$ называется \emph{расширением Галуа}, если оно нормально и
сепарабельно, т.\,е.\ минимальный многочлен над $F$ каждого элемента $\Om$ сепарабелен и раскладывается в $\Om$ на
линейные множители.
\end{definition}

Для конечных расширений это определение совпадает с прежним (\rf{3.9}). Бесконечные расширения Галуа возникают
естественно: таковы алгебраическое замыкание $\Fpb$ конечного поля, поле $\Qal$ всех алгебраических чисел,
сепарабельное замыкание $\Fsep$ любого поля $F$, поле $\Q(\mu_\infty)$, порождённое всеми корнями из единицы.

\paragraph{Как устроена группа автоморфизмов.} Пусть $\Om/F$~--- расширение Галуа и $F\subset M\subset\Om$, где
$M/F$~--- \emph{конечное} расширение Галуа. Каждый $\sigma\in\Aut(\Om/F)$ переводит $M$ в себя: $M$~--- поле разложения
некоторого сепарабельного многочлена $f\in F[X]$, а $\sigma$ переводит корни $f$ в корни $f$. Поэтому определено
ограничение $\sigma\mapsto\sigma|_M\colon\Aut(\Om/F)\to\Gal(M/F)$, а вместе с ним гомоморфизм
\begin{equation}\label{eq:Phi}
\Phi\colon\Aut(\Om/F)\to\prod_{M}\Gal(M/F),\qquad \sigma\mapsto(\sigma|_M)_M,
\end{equation}
где $M$ пробегает все конечные подрасширения Галуа в $\Om$.

\emph{$\Phi$ инъективен.} Каждый $\alpha\in\Om$ лежит в некотором конечном подрасширении Галуа: в поле разложения
$M_\alpha$ минимального многочлена $\alpha$ над $F$ (он сепарабелен, и все его корни лежат в $\Om$). Если
$\Phi(\sigma)=1$, то $\sigma$ тождествен на каждом $M_\alpha$, т.\,е.\ $\sigma(\alpha)=\alpha$ для всех $\alpha$, и
$\sigma=\id$.

Снабдим каждую конечную группу $\Gal(M/F)$ дискретной топологией. Она компактна и хаусдорфова, поэтому произведение
$\prod_M\Gal(M/F)$ с топологией произведения компактно (теорема Тихонова; здесь используется аксиома выбора) и
хаусдорфово. Мы увидим (\rf{7.8}, \rf{7.8a}), что образ $\Phi$ \emph{замкнут}; поэтому $\Aut(\Om/F)$ с перенесённой через $\Phi$
топологией~--- компактная группа. Эту группу с этой топологией обозначают $\Gal(\Om/F)$ и называют \emph{группой Галуа}
расширения $\Om/F$.

\begin{center}
\begin{tikzpicture}[>=stealth]
\node (A) at (0,0) {$\Aut(\Om/F)$};
\node (I) at (4.1,0) {$\Phi\bigl(\Aut(\Om/F)\bigr)$};
\node (P) at (8.4,0) {$\prod_{M}\Gal(M/F)$};
\draw[->,thick] (A) -- node[above]{$\Phi$} node[below,font=\scriptsize]{биекция} (I);
\node at ($(I.east)!0.5!(P.west)$) {$\subset$};
\draw[decorate,decoration={brace,mirror,amplitude=4pt},cnew] ($(A.west)+(0,-0.42)$) -- node[below=6pt,font=\footnotesize,align=center]{топология\\ переносится через $\Phi$} ($(A.east)+(0,-0.42)$);
\draw[decorate,decoration={brace,mirror,amplitude=4pt},cold] ($(I.west)+(0,-0.42)$) -- node[below=6pt,font=\footnotesize,align=center]{замкнуто\\ (7.8)} ($(I.east)+(0,-0.42)$);
\draw[decorate,decoration={brace,mirror,amplitude=4pt},cold] ($(P.west)+(0,-0.42)$) -- node[below=6pt,font=\footnotesize,align=center]{компактно: $\Gal(M/F)$ конечны\\ и дискретны; теорема Тихонова} ($(P.east)+(0,-0.42)$);
\end{tikzpicture}
\end{center}
\figcap{Замкнутое подмножество компакта компактно. Поэтому $\Aut(\Om/F)$ с перенесённой через $\Phi$ топологией~---
компактная группа.}

Главный результат главы (основная теорема \rf{7.13}): вся теория Галуа конечных расширений переносится на бесконечные, если
всюду заменить слово <<подгруппа>> словами <<\emph{замкнутая} подгруппа>>. Замкнутые подгруппы компактной группы сами компактны, так что
в бесконечной теории Галуа подгруппы~--- это компактные подгруппы $\Gal(\Om/F)$. Одно отличие следует отметить сразу: если
$\Om/F$ бесконечно, то $\Gal(\Om/F)$ всегда несчётна (замечание \rf{7.26b}), тогда как степень $[\Om:F]$ может быть счётной;
поэтому равенство $|\Gal(\Om/F)|=[\Om:F]$ перестаёт выполняться (хотя группа $\Gal(\Om/F)$
бесконечна тогда и только тогда, когда бесконечна степень $[\Om:F]$).

\begin{remark}
В этой главе свободно используется аксиома выбора (через лемму Цорна и теорему Тихонова). Без неё разумной теории не
получится: с аксиомами Цермело--Френкеля (ZF) совместимо существование алгебраического замыкания поля $\Q$, у
которого нет нетривиальных автоморфизмов.
\end{remark}
